\documentclass[11pt]{amsart}

\usepackage[T1]{fontenc}
\usepackage{amsmath,amssymb,amsthm,mathtools}
\usepackage[a4paper,left=22mm,right=22mm,top=21mm,bottom=23mm]{geometry}
\usepackage[hidelinks]{hyperref}
\usepackage{microtype}
\microtypesetup{expansion=false}

\newtheorem{theorem}{Theorem}[section]
\newtheorem{proposition}[theorem]{Proposition}
\newtheorem{lemma}[theorem]{Lemma}
\newtheorem{corollary}[theorem]{Corollary}

\newcommand{\Z}{\mathbb Z}
\newcommand{\Sym}{\operatorname{Sym}}
\newcommand{\gm}{\gamma}
\newcommand{\dlt}{\delta}
\newcommand{\aug}{\varpi}
\newcommand{\Rel}{\mathcal R}
\newcommand{\Alt}{\operatorname{Alt}}

\title[The fifth dimension subring]{The fifth dimension subring:\
an invariant quadratic proof}
\author{}
\date{August 2026}

\begin{document}

\begin{abstract}
We prove $\delta _5(L)\subseteq\gamma _4(L)$ for every Lie ring.
After a finite-data reduction, a class-three dimension witness is
converted into a cycle lifted along
$L/\gamma _3(L)\to L/\gamma _2(L)$.  A canonical alternating tensor
both evaluates that witness and proves that every lifted cycle has
value zero.  Quotients are written explicitly throughout; apart from
the free generators, auxiliary bases occur only in one tagged PBW
collection.
\end{abstract}

\maketitle

\section{Statement and reduction}

All tensor products, symmetric powers, and exterior powers are over
$\Z$.  For a Lie ring $L$, let $U(L)$ be its universal enveloping
algebra and $\aug(L)$ its augmentation ideal.  Define
\[
 \dlt_n(L)=\{x\in L: x\text{ maps into }\aug(L)^n\subseteq U(L)\},
 \qquad
 \gm_1(L)=L,\quad \gm_{n+1}(L)=[L,\gm_n(L)].
\]

\begin{theorem}\label{thm:main}
For every Lie ring $L$,
\[
                         \dlt_5(L)\subseteq\gm_4(L).
\]
\end{theorem}

We first prove the finitely generated class-three case directly; no
finite-order or primary hypothesis is needed.

\begin{proposition}[class-three case]\label{prop:class-three}
If $L$ is finitely generated as a Lie ring and $\gm_4(L)=0$, then
$\dlt_5(L)=0$.
\end{proposition}

\begin{lemma}[reduction to finite data]\label{lem:reduction}
Proposition~\ref{prop:class-three} implies
Theorem~\ref{thm:main}.
\end{lemma}

\begin{proof}
Given $L$, put $\bar L=L/\gm_4(L)$.  Functoriality gives
$\dlt_5(L)\to\dlt_5(\bar L)$, so it is enough to prove
$\dlt_5(\bar L)=0$.

The Lie ring $\bar L$ is the filtered union of its finitely generated
Lie subrings, and the left adjoint $U(-)$ commutes with this filtered
colimit.  If $a\in\dlt_5(\bar L)$, choose a finitely generated stage
$M_0$ containing $a$ and carrying a finite witness
$w\in\aug(M_0)^5$.  The images of $w$ and $a$ are equal in
$U(\bar L)=\varinjlim_MU(M)$, so they are equal in $U(M)$ at some
later finitely generated stage $M\supseteq M_0$.  Thus
$a\in\dlt_5(M)$.  Since $\gm_4(M)=0$,
Proposition~\ref{prop:class-three} gives $a=0$.
\end{proof}

For the rest of the proof, fix a finitely generated Lie ring $L$ with
$\gm_4(L)=0$.  Choose a finite free abelian group $X$ whose basis maps
to Lie generators of $L$, and let
\[
 F=\mathbb L(X)\twoheadrightarrow L,
 \qquad \Rel=\ker(F\to L).
\]
Write $F=\bigoplus_{j\geq1}F_j$ by bracket weight and put
$F_{\geq s}=\bigoplus_{j\geq s}F_j$.  Then
$F_1=X$, $F_2=\Lambda^2X$, and $F_{\geq4}\subseteq\Rel$.
The finite-rank group $F_1\oplus F_2\oplus F_3$ maps onto the
underlying additive group of $L$; in particular, all additive groups
used below are finitely generated.
An overline will always mean image in $L$.

\section{The invariant part of the proof}

\subsection{Two explicit resolutions}

The superscript on each resolution below records the denominator of
the quotient it resolves:
\begin{equation}\label{eq:resolutions}
 \begin{aligned}
 P^{(2)}_1&=\ker\bigl(X\longrightarrow L/\gm_2(L)\bigr)
       \xhookrightarrow{\ d_2\ }P^{(2)}_0=X,\\
 P^{(3)}_1&=\ker\bigl(X\oplus F_2\longrightarrow L/\gm_3(L)\bigr)
       \xhookrightarrow{\ d_3\ }P^{(3)}_0=X\oplus F_2.
 \end{aligned}
\end{equation}
The second displayed map sends $(x,y)$ to $\bar x+\bar y$ modulo
$\gm_3(L)$.  The quotient maps from $P^{(2)}_0$ and $P^{(3)}_0$ in
\eqref{eq:resolutions} are surjective: modulo $\gm_2(L)$ only
the chosen generators are needed, and modulo $\gm_3(L)$ only
generators and their weight-two brackets are needed.
\mbox{Each kernel in \eqref{eq:resolutions} is free}, being a subgroup of a finite free
abelian group.  Thus these are free resolutions of the two displayed
additive groups.

Projection $X\oplus F_2\to X$ gives a chain map resolving
$\pi:L/\gm_3(L)\to L/\gm_2(L)$:
\begin{equation}\label{eq:pi-chain}
 \pi_\bullet:P^{(3)}_\bullet\longrightarrow P^{(2)}_\bullet
\end{equation}
Indeed, if
$\bar x+\bar y\in\gm_3(L)$ with $y\in F_2$, then
$\bar x\in\gm_2(L)$, so projection sends $P^{(3)}_1$ into
$P^{(2)}_1$.

For a two-term free resolution $T=(T_1\xrightarrow tT_0)$, set
\begin{equation}\label{eq:Koszul}
 K_S(T)=\left[
 \Lambda^2T_1\xrightarrow{\kappa_2}T_1\otimes T_0
 \xrightarrow{\kappa_1}\Sym^2T_0\right],
\end{equation}
where
\[
 \kappa_2(a\wedge b)=a\otimes tb-b\otimes ta,
 \qquad
 \kappa_1(a\otimes x)=(ta)x.
\]
Write $Z_1K_S(T)=\ker\kappa_1$.
The quadratic Koszul theorem of K\"ock\footnote{B.~K\"ock,
\emph{Computing the homology of Koszul complexes}, Trans. Amer. Math.
Soc. \textbf{353} (2001), 3115--3147.} gives the natural identification
\begin{equation}\label{eq:derived-square}
 H_1K_S(T)=L_1\Sym^2(\operatorname{coker}t).
\end{equation}
Nothing below uses any other model of this derived functor.

\begin{lemma}[alternating tensors]\label{lem:Alt}
For every finitely generated abelian group $G$, the map
\[
 \Alt_G:\Lambda^2G\longrightarrow G\otimes G,
 \qquad x\wedge y\longmapsto x\otimes y-y\otimes x,
\]
is injective.  If $G$ is free, its image is
$\ker(G\otimes G\to\Sym^2G)$.
\end{lemma}

\begin{proof}
Write $G$ as a direct sum of cyclic groups $G_i$.  Since
$\Lambda^2G_i=0$, one has
$\Lambda^2G=\bigoplus_{i<j}G_i\otimes G_j$; projection of $\Alt_G$
to the $(i,j)$ tensor component is the identity on that summand.
This proves injectivity.  For free $G$, the second assertion is the
defining presentation of $\Sym^2G$.
\end{proof}

\subsection{The canonical quadratic invariant}

For $D$ in any of $\Lambda^2X$, $X\otimes F_2$, or $\Lambda^2F_2$,
write $[D]_F$ for bracket evaluation in $F$ and $[D]_L$ for its image
in $L$.  On $\Lambda^2X=F_2$, one simply has $[D]_F=D$.

Let $c\in Z_1K_S(P^{(2)})$.  Symmetric multiplication sends
$(d_2\otimes1)c$ to $\kappa_1c=0$.  Thus this tensor lies in
$\ker(X\otimes X\to\Sym^2X)$, and Lemma~\ref{lem:Alt} gives a unique
\begin{equation}\label{eq:alpha}
 \alpha(c)\in\Lambda^2X,
 \qquad
 \Alt_X\alpha(c)=(d_2\otimes1)c.
\end{equation}

\begin{lemma}[canonical invariant]\label{lem:eta}
The bracket value $[\alpha(c)]_L$ belongs to $\gm_3(L)$, and
\begin{equation}\label{eq:eta}
 \eta_L:L_1\Sym^2\bigl(L/\gm_2(L)\bigr)\longrightarrow\gm_3(L),
 \qquad \eta_L([c])=-[\alpha(c)]_L,
\end{equation}
is a well-defined homomorphism, natural among finitely generated Lie
rings with $\gm_4=0$.
\end{lemma}

\begin{proof}
Apply $X\to L/\gm_2(L)$ to both factors of \eqref{eq:alpha}.  Its
right side becomes zero because $d_2P^{(2)}_1$ maps to zero.  By
naturality and injectivity of $\Alt$,
$\alpha(c)$ has zero image in
$\Lambda^2(L/\gm_2(L))$.  The bracket map
\[
 \Lambda^2X\longrightarrow\gm_2(L)/\gm_3(L)
\]
factors through $\Lambda^2(L/\gm_2(L))$: changing either entry by an
element of $\gm_2(L)$ changes its bracket by an element of
$\gm_3(L)$.  Thus $[\alpha(c)]_L\in\gm_3(L)$.

If $u,v\in P^{(2)}_1$, then directly from the definitions
\[
 \alpha\bigl(\kappa_2(u\wedge v)\bigr)=d_2u\wedge d_2v.
\]
Both entries map into $\gm_2(L)$, so their bracket lies in
$[\gm_2(L),\gm_2(L)]\subseteq\gm_4(L)=0$.  Formula \eqref{eq:eta}
therefore kills boundaries.

Call two presentations compatible when a homomorphism of their
generator groups induces a map of free Lie rings commuting with the
maps to the given Lie rings.  Formula \eqref{eq:alpha} is natural for
such maps.  For two presentations of the same $L$, use the one on the
direct sum of their generator groups as a common target.  The two
inclusions are compatible and induce the identity on
$L_1\Sym^2(L/\gm_2(L))$, so they give the same $\eta_L$.  For a map
$L\to L'$, add to a presentation of $L'$ one generator representing
the image of each chosen generator of $L$.  This gives a compatible
map into the enlarged target presentation.  Comparing along the
inclusion of the original target presentation into the enlarged one
proves naturality for the original map.
\end{proof}

\begin{lemma}[calculus of a lifted cycle]\label{lem:lifted-cycle}
Let $\chi\in Z_1K_S(P^{(3)})$ and let
\[
 c=(\pi_1\otimes\pi_0)\chi\in Z_1K_S(P^{(2)}).
\]
There is a unique
\begin{equation}\label{eq:beta}
 \beta(\chi)\in\Lambda^2(X\oplus F_2),
 \qquad
 \Alt_{X\oplus F_2}\beta(\chi)=(d_3\otimes1)\chi.
\end{equation}
Write, according to the two summands $F_1=X$ and $F_2$,
\begin{equation}\label{eq:beta-split}
 \beta(\chi)=\beta_{11}+\beta_{12}+\beta_{22}
 \in\Lambda^2X\oplus(X\otimes F_2)\oplus\Lambda^2F_2.
\end{equation}
Then
\begin{align}
 \beta_{11}&=\alpha(c),                                      \label{eq:beta-alpha}\\
 [\beta_{12}]_L&=-2[\beta_{11}]_L,                           \label{eq:cross}\\
 [\beta_{12}]_L&=[\beta_{22}]_L=0.                          \label{eq:mixed-zero}
\end{align}
The cross identity \eqref{eq:cross} evaluates the PBW witness; the
mixed vanishing \eqref{eq:mixed-zero} is used in the invariant
vanishing argument.
\end{lemma}

\begin{proof}
Symmetric multiplication sends $(d_3\otimes1)\chi$ to
$\kappa_1\chi=0$.  Existence and uniqueness in \eqref{eq:beta}
therefore follow from Lemma~\ref{lem:Alt}.  Projecting
\eqref{eq:beta} to $X\otimes X$ gives \eqref{eq:alpha}, hence
\eqref{eq:beta-alpha}.

For the remaining identities, write
\[
 \chi=\sum_i r_i\otimes(u_i,v_i),
 \qquad d_3r_i=(x_i,y_i),
 \quad x_i,u_i\in X,\quad y_i,v_i\in F_2.
\]
The $X\otimes X$, $F_2\otimes X$, and $X\otimes F_2$ components of
\eqref{eq:beta} are respectively
\begin{equation}\label{eq:three-components}
 \sum_i x_i\otimes u_i=\Alt_X\beta_{11},
 \qquad
 \sum_i y_i\otimes u_i=-\tau(\beta_{12}),
 \qquad
 \sum_i x_i\otimes v_i=\beta_{12},
\end{equation}
where $\tau$ exchanges the tensor factors.  Since $r_i\in P^{(3)}_1$,
the element $\bar x_i+\bar y_i$ lies in $\gm_3(L)$, which is central
because $[L,\gm_3(L)]=\gm_4(L)=0$; hence
$[\bar y_i,\bar u_i]=-[\bar x_i,\bar u_i]$.  Applying the bracket
to the first two equations in \eqref{eq:three-components} gives
\[
 \sum_i[\bar x_i,\bar u_i]=2[\beta_{11}]_L,
 \qquad
 \sum_i[\bar y_i,\bar u_i]=[\beta_{12}]_L.
\]
This proves \eqref{eq:cross}, including its sign and factor $2$.

Moreover, $\bar x_i\in\gm_2(L)$ because
$\bar x_i+\bar y_i\in\gm_3(L)$ and $\bar y_i\in\gm_2(L)$.
Thus both entries in the last equation of
\eqref{eq:three-components} map into $\gm_2(L)$, proving the first
equality in \eqref{eq:mixed-zero}; the second is immediate for
$\beta_{22}\in\Lambda^2F_2$.  Here both brackets vanish because
$[\gm_2(L),\gm_2(L)]\subseteq\gm_4(L)=0$.
\end{proof}

\begin{corollary}[invariant vanishing]\label{cor:vanishing}
For $\pi:L/\gm_3(L)\to L/\gm_2(L)$,
\begin{equation}\label{eq:vanishing}
 \eta_L\circ L_1\Sym^2(\pi)=0.
\end{equation}
\end{corollary}

\begin{proof}
Represent a class in $L_1\Sym^2(L/\gm_3(L))$ by a cycle $\chi$ and
map \eqref{eq:beta} to
$(L/\gm_3(L))\otimes(L/\gm_3(L))$.  Its right side is zero because
$d_3P^{(3)}_1$ maps to zero.  Lemma~\ref{lem:Alt} says that the image
of $\beta(\chi)$ in $\Lambda^2(L/\gm_3(L))$ is zero.  Since
$\gm_3(L)$ is central, the bracket of lifts gives a well-defined map
from this exterior square to $\gm_2(L)$.  Applying it and using
\eqref{eq:mixed-zero} gives
\[
 0=[\beta_{11}]_L+[\beta_{12}]_L+[\beta_{22}]_L=[\beta_{11}]_L.
\]
Here the exterior bracket sends $x\wedge y$ to $[x,y]$, so no factor
$2$ occurs.  Now
\eqref{eq:beta-alpha} and \eqref{eq:eta} prove \eqref{eq:vanishing}.
\end{proof}

\section{The one PBW calculation}

For $u\in U(F)=T(X)$, let $u_j$ be its homogeneous tensor-weight
$j$ component.  Truncating a relation gives
\begin{equation}\label{eq:lambda}
 \lambda:\Rel\longrightarrow P^{(3)}_1,
 \qquad \lambda(\rho)=(\rho_1,\rho_2).
\end{equation}

\begin{lemma}[relation truncation]\label{lem:lambda}
The map \eqref{eq:lambda} is well defined and surjective.
\end{lemma}

\begin{proof}
For $\rho\in\Rel$, the image of $\rho_1+\rho_2$ in $L$ is the
negative of the image of $\rho_{\geq3}$ and therefore belongs to
$\gm_3(L)$.  Thus $\lambda(\rho)\in P^{(3)}_1$.  Conversely, if
$(x,y)\in P^{(3)}_1$, then $\bar x+\bar y\in\gm_3(L)$.  Since
$F_{\geq3}=\gm_3(F)$ maps onto $\gm_3(L)$ and $F_{\geq4}$ maps to
zero, the map $F_3\to\gm_3(L)$ is onto.  Choose $z\in F_3$ with
$\bar z=-\bar x-\bar y$.  Then $x+y+z\in\Rel$ and has truncation
$(x,y)$.
\end{proof}

Let $\operatorname{mult}:U(F)\otimes U(F)\to U(F)$ denote
multiplication; in particular,
$\operatorname{mult}(\rho\otimes h)=\rho h$.  Apart from the chosen
free generators, the following lemma is the only place where
auxiliary bases are chosen.

\begin{lemma}[placed PBW collection]\label{lem:placed}
Suppose $\theta\in\Rel U(F)$ and $\theta_j\in F_j$ for
$1\leq j\leq4$.  There exist
\[
 \sigma\in\Rel,
 \qquad \chi_{\mathrm{tag}}\in\Rel\otimes(X\oplus F_2),
 \qquad E_{\mathrm{rem}}\in U(F)
\]
such that, modulo tensor weights at least five,
\begin{equation}\label{eq:placed-decomposition}
 \theta=\sigma+\operatorname{mult}(\chi_{\mathrm{tag}})+E_{\mathrm{rem}}.
\end{equation}
Moreover, the remainder $E_{\mathrm{rem}}$ has neither
\begin{enumerate}
 \item a one-factor PBW term in weights $1,2,3$, nor
 \item a two-factor PBW term in
 \begin{equation}\label{eq:low-block}
 \Sym^2X\oplus(X\otimes F_2)\oplus\Sym^2F_2
       =\Sym^2(X\oplus F_2).
 \end{equation}
\end{enumerate}
\end{lemma}

\begin{proof}
\emph{Relation rows.}
For $1\leq s\leq4$, put
\[
 D_s=\operatorname{im}(\Rel\cap F_{\geq s}\longrightarrow F_s).
\]
By Smith normal form, choose a basis of $F_s$ in which a basis of
$D_s$ consists of integer multiples of distinct basis elements.  Lift
that basis of $D_s$ to relations and call the lifts rows.  A row of
least weight $s$ has the corresponding multiple as its first
homogeneous term; call it the Smith head.  Successively
subtracting rows to remove weights $1,2,3,4$ shows that the rows
generate $\Rel$ modulo $F_{\geq5}$, integrally.  Multiplication by
external factors cannot lower the weight of the residual.

The truncations of the rows of least weights one and two form a
$\Z$-basis of $P^{(3)}_1$.  They span because Lemma~\ref{lem:lambda}
lifts any element of $P^{(3)}_1$ to a relation, after which the
weight-one and weight-two components can be removed successively.
They are independent because the weight-one components first force
all coefficients of weight-one rows to vanish, and the weight-two
components then do the same for the remaining rows.

\emph{Terminating tagged collection.}
Order the chosen homogeneous bases of $F$ first by weight and use the
corresponding Poincar\'e--Birkhoff--Witt (PBW) basis.  A PBW factor
means one homogeneous Lie-basis factor; call each marked summand a
packet.  Express $\theta$, modulo
weight five, as a sum of packets and collect while retaining the mark,
using
\begin{equation}\label{eq:collection-identities}
 h\rho=\rho h+[h,\rho],
 \qquad \rho h=h\rho+[\rho,h],
 \qquad h'h=hh'-[h,h']\quad(h<h').
\end{equation}
Every commutator with a row is again a relation and is immediately
re-expanded in the fixed rows.  Regard that normalization as part of
the rewrite that created the relation.  If $s$ is the least row
weight of the parent packet, $r$ its number of external factors, and
$I$ its inversion number after replacing the mark by its Smith head,
every normalized output packet is smaller lexicographically in
\[
                              (4-s,r,I).
\]
Indeed a row remainder raises $s$, a commutator output has fewer
external factors, and an ordered principal output has fewer
inversions.  Initial relations are already fixed rows.  Terms
reaching weight five are discarded, so the process terminates.
Fixing the leftmost available rewrite makes it unambiguous; no
confluence claim is required.

For reference, all possible external factor-weight sequences below
weight five are
\begin{equation}\label{eq:source-table}
\begin{array}{c|l}
\text{least row weight }s&\text{weights of the external factors}\\ \hline
1&(),(1),(2),(1,1),(3),(1,2),(1,1,1)\\
2&(),(1),(2),(1,1)\\
3&(),(1)\\
4&().
\end{array}
\end{equation}
Here $()$ is the empty sequence.  The table is simply the list of
partitions of an external weight at most $4-s$, with factors in PBW
order.

\pagebreak
\emph{The frozen packets.}
Select every terminal packet having exactly one external factor
$h\in X\oplus F_2$ and a row of least weight one or two.  Orient it
as $\rho h$ by \eqref{eq:collection-identities}; put the whole
correction relation, with its sign, into $\sigma$, and leave the full product
$\rho h$ unexpanded.  With the actual integral coefficients of these
packets, put
\[
                         \chi_{\mathrm{tag}}=\sum n\rho\otimes h.
\]
Put all remaining packets with no external factor (scalar relations)
into $\sigma$ and fully collect
the unselected packets.  Their sum is $E_{\mathrm{rem}}$, giving
\eqref{eq:placed-decomposition}.  Because the relation tags were
retained before any PBW coefficients were combined, this construction
uses no saturation or divisibility inference.

\emph{Verification of the two exclusions.}
Every newly created packet with no external factor goes into $\sigma$,
so $E_{\mathrm{rem}}$ has no one-factor term; every newly created
eligible one-external packet is frozen as above.  An unselected packet
with one external factor and total weight below five has two
head-factor weights $1$ and $3$.  Its symbol lies outside
\eqref{eq:low-block}, and every higher row term has weight at least
five.  A packet with at least two external factors has at least three
PBW factors in its head term.  If
its row has least weight $s$ and its $r\geq2$ external factors have
total weight $e\geq r$, a higher row term starts in weight
$s+1+e$.  This is below five only for $s=1$ and $e=r=2$.  The sole
tail to check is therefore $Yxx'$ with
$Y\in F_2$ and two already ordered elements $x\leq x'$ of $X$:
\[
 Yxx'=xx'Y+x[Y,x']+[Y,x]x'.
\]
Further ordering yields only three-factor terms, two-factor terms of
type $X F_3$, and a one-factor term in $F_4$; it never yields a term
in $\Sym^2F_2$.  The $F_4$ term is a scalar relation and is absorbed
into $\sigma$.  If the two external $X$'s had earlier been inverted,
their commutator would already have produced a selected one-external
packet $\rho F_2$.  This proves both asserted exclusions and exhausts
the table \eqref{eq:source-table}.
\end{proof}

\begin{lemma}[homological extraction]\label{lem:extraction}
Under the hypotheses of Lemma~\ref{lem:placed}, there are
$\sigma\in\Rel$ and $\chi\in Z_1K_S(P^{(3)})$ such that, for
$\beta(\chi)$ from \eqref{eq:beta},
\begin{equation}\label{eq:ledger}
 \theta_1=\sigma_1,
 \qquad
 \theta_2=\sigma_2+[\beta_{11}]_F,
 \qquad
 \theta_3=\sigma_3+[\beta_{12}]_F.
\end{equation}
Moreover the projection of $\chi$ to $K_S(P^{(2)})$ is a cycle and
$\beta_{11}$ is its tensor $\alpha$ from \eqref{eq:alpha}.
\end{lemma}

\begin{proof}
Take the output of Lemma~\ref{lem:placed} and set
\[
                     \chi=(\lambda\otimes1)\chi_{\mathrm{tag}}.
\]
Under the PBW two-factor-symbol identification, the low-block symbol
of $(\rho_1+\rho_2)h$ is
$\kappa_1((\lambda\rho)\otimes h)$.  A tail
$\rho_{\geq3}X$ has type $F_{\geq3}X$, outside
\eqref{eq:low-block}, while $\rho_{\geq3}F_2$ has weight at least
five.  Thus the component of
$\operatorname{mult}(\chi_{\mathrm{tag}})$ in
\eqref{eq:low-block} is exactly $\kappa_1\chi$.  The corresponding
components of $\theta$ and $\sigma$ vanish because their components
through weight four are Lie elements, while that of $E_{\mathrm{rem}}$
vanishes by Lemma~\ref{lem:placed}(2).  Hence $\kappa_1\chi=0$, so
$\chi$ is a cycle.

The difference $\rho-(\rho_1+\rho_2)$ starts in weight three.  Any
one-factor term formed from this tail is a commutator with an external
factor of weight at least one, hence has weight at least four.  By
Lemma~\ref{lem:placed}(1), the one-factor components of
\eqref{eq:placed-decomposition} in weights at most three are thus
those of
\[
 \sigma+\operatorname{mult}\bigl((d_3\otimes1)\chi\bigr)
 =\sigma+\operatorname{mult}\bigl(\Alt_{X\oplus F_2}\beta(\chi)\bigr).
\]
Multiplication sends an alternating tensor $u\otimes v-v\otimes u$
to $[u,v]$.  The three components in \eqref{eq:beta-split} therefore
give respectively the one-factor terms
$[\beta_{11}]_F$, $[\beta_{12}]_F$, and $[\beta_{22}]_F$ in weights
$2,3,4$.  This proves \eqref{eq:ledger}.  Finally, the chain map
$\pi_\bullet$ sends $\chi$ to a cycle, and projection of
\eqref{eq:beta} to $X\otimes X$ proves the last assertion.
\end{proof}

\section{A dimension witness and the conclusion}

\begin{theorem}[normal form]\label{thm:normal}
For every $a\in\dlt_5(L)$ there is
\[
 \xi\in L_1\Sym^2\bigl(L/\gm_3(L)\bigr)
\]
such that
\begin{equation}\label{eq:normal}
 a=\eta_L\bigl(L_1\Sym^2(\pi)(\xi)\bigr),
 \qquad \pi:L/\gm_3(L)\longrightarrow L/\gm_2(L).
\end{equation}
\end{theorem}

\begin{proof}
Choose
$f=f_1+f_2+f_3\in F_1\oplus F_2\oplus F_3$ mapping to $a$; this is
possible because $F_{\geq4}\subseteq\Rel$.  The surjection
$U(F)\to U(L)$ maps $\aug(F)^5$ onto $\aug(L)^5$, so choose
$t\in\aug(F)^5$ with the same image as $f$, and put
\[
                              \theta=f-t.
\]
The kernel of $U(F)\to U(L)$ is the two-sided ideal generated by
$\Rel$.  Since $x\rho=\rho x+[x,\rho]$ and
$[F,\Rel]\subseteq\Rel$, induction on left-word length shows that
this kernel is the right ideal $\Rel U(F)$.
Thus $\theta\in\Rel U(F)$.  Since $U(F)=T(X)$ and
$t\in\aug(F)^5$, one has
\[
                 \theta_j=f_j\ (j=1,2,3),\qquad \theta_4=0.
\]

Apply Lemma~\ref{lem:extraction}.  Summing its three ledger equations
and mapping to $L$ gives
\begin{equation}\label{eq:witness-evaluation}
 a=\overline{\sigma_1+\sigma_2+\sigma_3}
      +[\beta_{11}]_L+[\beta_{12}]_L
   =[\beta_{11}]_L+[\beta_{12}]_L.
\end{equation}
The second equality holds because $\sigma\in\Rel$, while all its
components of weight at least four map into $\gm_4(L)=0$; explicitly,
$0=\bar\sigma=\overline{\sigma_1+\sigma_2+\sigma_3}$.
Equation \eqref{eq:cross} now yields
\[
                 a=-[\beta_{11}]_L
                   =-[\alpha(c)]_L=\eta_L([c]),
\]
where $c=(\pi_1\otimes\pi_0)\chi$.  Under the identifications
\eqref{eq:derived-square}, $[c]$ is the image of
$\xi=[\chi]\in L_1\Sym^2(L/\gm_3(L))$.  This is \eqref{eq:normal}.
\end{proof}

\begin{proof}[Proof of Proposition~\ref{prop:class-three}]
For $a\in\dlt_5(L)$, Theorem~\ref{thm:normal} gives
$a=(\eta_L\circ L_1\Sym^2(\pi))(\xi)$ for some $\xi$.
Corollary~\ref{cor:vanishing} makes this value zero.
\end{proof}

\begin{proof}[Proof of Theorem~\ref{thm:main}]
Apply Lemma~\ref{lem:reduction} and
Proposition~\ref{prop:class-three}.
\end{proof}

\end{document}
